\chapter{Literature Review and Theoretical Foundations}
\label{ch:literature_review}

Recognizing handwritten mathematics has developed over several decades, moving from rigid, rule-based systems to modern deep neural networks. Even so, taking models that work on single isolated equations and getting them to transcribe messy, full-page university exam booklets is a massive leap. In this chapter, we trace the history of Handwritten Mathematical Expression Recognition (HMER), walk through the math behind attention-based sequence decoders, explain our early experiment with GAN-based font normalization and why it failed (Pivot 1), and look at why single-formula models struggle on full exam pages. Finally, we review document layout analysis and recent multimodal vision-language architectures.

\section{How Math Recognition Evolved}
\label{sec:chronological_evolution}

Early systems for reading math relied on dividing the problem into three separate steps \cite{smith2007overview}:
\begin{enumerate}
    \item \textbf{Symbol Segmentation:} The system binarized the image and broke it into individual character components using projection profiles, connected components, or contour analysis.
    \item \textbf{Isolated Symbol Recognition:} Each detected component was sent to a classifier (such as an SVM or a shallow CNN) to predict the character.
    \item \textbf{Grammar and Structural Parsing:} Once symbols were identified, a 2D stochastic context-free grammar or minimum spanning tree tried to figure out their spatial relationships (like superscripts, subscripts, fractions, or matrix rows).
\end{enumerate}

\subsection{Why Classical Three-Stage Pipelines Break Down}
While this three-stage design makes intuitive sense, it suffered from a major flaw: \textbf{error propagation}. In handwritten math, you cannot reliably tell what a stroke is without knowing the math around it. A single horizontal line could be a minus sign ($-$), a fraction bar ($\frac{a}{b}$), the crossbar of a `$t$', or an accent bar ($\bar{x}$). If the initial segmentation step grouped strokes incorrectly, the downstream grammar parser had no way to recover. On top of that, 2D grammar parsing took exponential time as equations grew longer, making it unusable for multi-line exam solutions.

\section{The Shift to Sequence-to-Sequence Attention}
\label{sec:seq2seq_revolution}

The field changed dramatically when researchers adopted end-to-end differentiable sequence models, borrowing ideas from neural machine translation \cite{vaswani2017attention}. In 2017, Deng et al. \cite{deng2017image} introduced the Im2Latex formulation, showing that formula recognition could be treated as translating an image directly into \LaTeX\ markup. Around the same time, Zhang et al. \cite{zhang2017watch} published the \textbf{Watch, Attend and Parse (WAP)} model, which became the baseline architecture for modern HMER.

\subsection{Mathematical Formulation of WAP}
WAP threw away explicit symbol segmentation. Instead, it paired a deep convolutional visual encoder with an attention-based recurrent decoder.

Given a normalized grayscale image $\mathbf{I} \in \mathbb{R}^{H \times W \times 1}$, a convolutional feature extractor (like DenseNet \cite{huang2017densely} or ResNet \cite{he2016deep}) maps the pixels into a 2D feature grid:
\begin{equation}
\mathbf{F} = \text{Encoder}(\mathbf{I}) \in \mathbb{R}^{H' \times W' \times D}
\end{equation}
where $H' = \lfloor H / 16 \rfloor$, $W' = \lfloor W / 16 \rfloor$, and $D$ is the feature channel depth. We flatten this feature grid into a set of visual tokens $\mathbf{V} = \{\mathbf{v}_1, \dots, \mathbf{v}_L\}$, where $L = H' \times W'$.

At each decoding step $t$, the decoder uses its hidden state $\mathbf{s}_t$ to compute spatial attention weights $\alpha_{t, i}$ over all visual tokens:
\begin{equation}
\alpha_{t, i} = \frac{\exp(e_{t, i})}{\sum_{j=1}^L \exp(e_{t, j})}, \quad \mathbf{c}_t = \sum_{i=1}^L \alpha_{t, i} \mathbf{v}_i
\end{equation}
Here, $e_{t, i}$ scores how well decoder state $\mathbf{s}_t$ matches visual feature $\mathbf{v}_i$, and $\mathbf{c}_t$ is the weighted context vector. The decoder then predicts the next \LaTeX\ token from vocabulary $\mathcal{V}_{\text{LaTeX}}$:
\begin{equation}
P(y_t \mid y_{<t}, \mathbf{I}) = \text{softmax}\left( \mathbf{W}_o [\mathbf{s}_t; \mathbf{c}_t] \right)
\end{equation}
By training the whole network end-to-end with cross-entropy loss, WAP learned to spot symbols and understand their 2D layout at the same time, without needing anyone to label individual strokes.

\section{The Attention Coverage Problem and Modern HMER}
\label{sec:coverage_deficit}

Even though WAP was a big improvement over classical pipelines, standard visual attention has a known weak spot when dealing with dense math: it has no memory of where it has already looked. Researchers call this the \textbf{Attention Coverage Deficit}.

\subsection{Under-Parsing and Over-Parsing}
Because standard attention calculates weights from scratch at every step, the decoder can easily lose track of its position on the page. In handwritten math, this leads to two opposite errors:
\begin{enumerate}
    \item \textbf{Under-Parsing (Skipping Small Symbols):} In mathematics, tiny symbols carry huge meaning. A minus sign ($-$), a decimal point ($.$), a prime mark ($f'$), or an exponent ($x^2$) takes up only a few pixels. In downsampled feature maps, their signal is faint. The attention mechanism often skips over them entirely, turning $x^2 - 1$ into $x 1$.
    \item \textbf{Over-Parsing (Infinite Loops):} When an equation contains repeating patterns---such as a row of zeros in Gaussian elimination ($0, 0, 0, 0$) or repeated variables---the decoder gets confused. It attends back to an earlier identical stroke and starts repeating the same token in a loop until hitting the maximum sequence length.
\end{enumerate}

\subsection{CoMER: Tracking Attention with Coverage Maps}
To tackle this issue, Zhao et al. \cite{zhao2022comer} introduced \textbf{CoMER (ECCV 2022)}. CoMER replaced the RNN decoder with a Transformer and added an \textbf{Attention Refinement Module (ARM)}.

CoMER keeps a running tally of where attention has been paid across previous steps by maintaining a cumulative coverage map $\mathbf{C}_t \in \mathbb{R}^{H' \times W'}$:
\begin{equation}
\mathbf{C}_t = \sum_{\tau=1}^{t-1} \boldsymbol{\alpha}_\tau
\end{equation}
This coverage map shows how much attention each pixel location has received so far. The ARM runs small convolutions over $\mathbf{C}_t$ to penalize areas that have already been read. If a region has already been transcribed ($\mathbf{C}_{t, i} \gg 1$), its attention score is pushed down toward zero, forcing the decoder to focus on unparsed ink.

Other recent models build on similar ideas: CAN (ECCV 2022) \cite{li2022can} adds an auxiliary task that counts symbols in the image, and PosFormer (ECCV 2024) \cite{zhao2024posformer} adds explicit 2D spatial coordinate embeddings into self-attention. Both show that explicitly tracking 2D positions is crucial for complex math.

\section{Our Early Experiment (Pivot 1): GAN Font Normalization}
\label{sec:pivot1_gan}

When we began this project in July and August 2026, we wondered if we could simplify the problem by turning messy handwriting into something existing OCR engines already handle well: \textbf{printed textbook math}.

Tools like \texttt{pix2tex} (LaTeX-OCR) \cite{blecher2023nougat} work well on crisp, computer-generated Computer Modern fonts. This led us to test our first idea: a \textbf{Two-Stage Generative Font Normalization Pipeline}.

\subsection{How the GAN Pipeline Worked}
We split the task into two neural network stages:
\begin{enumerate}
    \item \textbf{Stage 1 (Handwriting to Printed Image):} An image-to-image translation network (a U-Net generator with a PatchGAN discriminator \cite{isola2017image}) that took a handwritten equation crop $\mathbf{I}_{\text{hw}}$ and tried to turn it into a synthetic image $\mathbf{I}_{\text{font}}$ that looked like clean Computer Modern print.
    \item \textbf{Stage 2 (Printed Image to \LaTeX):} A Vision Transformer (ViT) that read the cleaned synthetic image and generated \LaTeX\ code.
\end{enumerate}

We trained Stage 1 with a combined loss function:
\begin{equation}
\mathcal{L}_G = \lambda_1 \|\mathbf{I}_{\text{font}} - \hat{\mathbf{I}}_{\text{font}}\|_1 + \lambda_{\text{perc}} \mathcal{L}_{\text{perceptual}} + \lambda_{\text{adv}} \mathcal{L}_{\text{adv}}
\end{equation}

\subsection{Why Generating Synthetic Pixels Failed}
At first glance, the generated images looked surprisingly clean and sharp. But when we inspected them carefully on real student handwriting, we found severe problems:
\begin{enumerate}
    \item \textbf{Plausibility over Accuracy:} A GAN is trained to fool its discriminator into thinking an image looks real. It does not care about mathematical correctness. To make a messy cursive stroke look like a neat printed letter, the generator would invent phantom serifs. It repeatedly turned cursive `$l$' into the numeral `$1$', or curled a `$v$' into an `$r$'.
    \item \textbf{Erasing Minus Signs and Exponents:} A minus sign or a decimal point covers very few pixels. In an $\mathcal{L}_1$ pixel loss, omitting a tiny dot incurs almost zero penalty compared to misdrawing a large capital letter. As a result, the GAN routinely smoothed over or erased minus signs and exponents entirely ($x^2 \to x$).
    \item \textbf{Smudging Matrix Delimiters:} When a student wrote a multi-row matrix with a slight slant, the GAN merged adjacent matrix columns into blurry blobs that the Stage 2 OCR model could not read.
\end{enumerate}

\subsection{Pivot 1: Ditching Intermediate Image Generation}
These experiments taught us a valuable lesson: \textbf{generating an intermediate synthetic image adds unnecessary noise and loses fine mathematical details}. Once a GAN hallucinates away a minus sign, no downstream model can recover it.

In August 2026, we made our \textbf{First Strategic Pivot (Pivot 1)}: we scrapped the image-to-image translation approach entirely. From then on, all our architectures mapped visual features directly into discrete \LaTeX\ tokens using end-to-end differentiable sequence models.

\section{Why Single-Formula Models Struggle on Real Exam Pages}
\label{sec:generalization_gap}

Even though models like CoMER score high on standard benchmarks, there is a big gap between academic benchmarks and real university exam booklets:
\begin{enumerate}
    \item \textbf{Tightly Cropped Formulas vs. Full Pages:} Benchmarks like CROHME \cite{mouchere2016icfhr, mahdavi2019crohme} and FERMAT \cite{soni2024fermat} contain only isolated, tightly cropped equations (like $\frac{a}{b} + c = d$) centered on clean backgrounds.
    \item \textbf{English Proofs and Mixed Explanations:} In university courses like MA102, students do not just write standalone formulas. They write continuous mathematical proofs mixing cursive English sentences (*"Let $W$ be a subspace of $V$..."*) with inline equations and matrix reductions. Models trained only on CROHME have never seen English words, so they try to parse cursive prose as mangled math symbols.
    \item \textbf{Downsampling Kills Details:} A full exam page scanned at 300 DPI is $2338 \times 1653$ pixels and holds 25 to 35 lines of writing. If you shrink that whole page down to the standard input size of an HMER model ($256 \times 512$ pixels), delicate pen strokes shrink to less than a pixel, making them impossible to read.
\end{enumerate}

This showed us that \textbf{an isolated HMER model alone cannot transcribe a full exam page}. We needed either a layout analysis engine to cut the page into clean lines and classify text versus math, or a native high-resolution vision-language model that could process the entire page at once.

\section{Document Layout Analysis on Unruled Exam Scripts}
\label{sec:layout_analysis}

To transcribe an entire page using modular models, the system must first find and segment the lines of writing. Traditional Document Layout Analysis (DLA) uses geometric heuristics:
\begin{enumerate}
    \item \textbf{Projection Profiles:} Summing pixels horizontally to find empty white spaces between lines works well for printed books. But on unruled paper, student writing drifts diagonally. The end of one line often overlaps vertically with the start of the next line, wiping out the white-space gaps.
    \item \textbf{Run-Length Smearing (RLSA):} Smearing black pixels along the horizontal axis to form text blocks fails when tall bracket delimiters (like $[ \dots ]$ for matrices) span four or five lines. The smearing algorithm merges the entire calculation into one giant unreadable block.
    \item \textbf{Interlocking Ascenders and Descenders:} Handwritten letters frequently poke into adjacent lines. The descender of a `$y$' or `$g$' from line 1 reaches down into the ascender of a `$d$' or `$\lambda$' on line 2 without actually touching it. A simple rectangular bounding box will slice these letters in half.
\end{enumerate}

\subsection{1D Vertical IoU Non-Maximum Suppression}
To stop letters from being cut in half without training a heavy pixel-segmentation model, we designed a \textbf{1D Vertical IoU Non-Maximum Suppression} algorithm.

Instead of working with 2D bounding boxes, we project each candidate line slice onto the vertical $y$-axis as an interval $J_k = [y_{k, \text{top}}, y_{k, \text{bottom}}]$. If two candidate slices $J_a$ and $J_b$ have a high vertical overlap:
\begin{equation}
\text{IoU}_{1D}(J_a, J_b) = \frac{\text{Length}(J_a \cap J_b)}{\text{Length}(J_a \cup J_b)} \ge \theta_{\text{IoU}} = 0.35
\end{equation}
the algorithm merges them into a single line container. This keeps interlocking letters together in the same line box so the recognition models can read them properly.

\section{Vision Transformers for Document OCR}
\label{sec:vision_transformers}

Recently, Transformer-based vision models have become the standard tool for document processing.

\subsection{Microsoft TrOCR and Vocabulary Drift}
Li et al. \cite{li2023trocr} introduced \textbf{TrOCR (AAAI 2023)}, which connects a Vision Transformer (ViT) encoder to an autoregressive Transformer text decoder. TrOCR works well on typical handwriting, but when we tested it on MA102 exam scripts, we uncovered a surprising issue: \textbf{historical vocabulary drift}.

The public TrOCR handwriting checkpoints were pretrained on the IAM database and historical manuscripts from the 18th and 19th centuries. When TrOCR encounters rapid cursive handwriting containing mathematical abbreviations (like $M_{4 \times 4}(\mathbb{R})$ or subspace conditions), its visual encoder becomes unsure. Instead of staying grounded in math, the decoder falls back on its pretraining data, suddenly producing archaic legal and political phrases like \emph{``House of Representatives''} or \emph{``Chamberlain v.''} right in the middle of a linear algebra proof. This proved that off-the-shelf text OCR models cannot be used on university math without careful domain constraints.

\subsection{Meta Nougat}
Blecher et al. \cite{blecher2023nougat} introduced \textbf{Nougat (Meta AI, 2023)}, pairing a Swin Transformer encoder with an mBART decoder to transcribe academic PDFs into Markdown/\LaTeX. Nougat works great on printed papers, but it fails completely on handwritten scripts. Because it was trained exclusively on digital vector typography, the variable stroke widths, uneven ink, and curved baselines of handwritten exams cause it to lose track and fall into repetitive generation loops.

\section{Multimodal Vision-Language Models (Qwen2.5-VL)}
\label{sec:vlm_frontier}

The latest development in document AI is the rise of native \textbf{Multimodal Large Language Models (MLLMs)} such as \textbf{Qwen2.5-VL} \cite{wang2024qwen2vl}, which can take an entire high-resolution image as input.

\subsection{Dynamic Resolution and 2D-RoPE}
Earlier vision models forced every image into a fixed square size (like $224 \times 224$ or $448 \times 448$), stretching or squishing non-square pages and distorting tiny mathematical symbols. Qwen2.5-VL fixes this with a \textbf{Dynamic Resolution ViT} that splits images into patches while preserving their original aspect ratio.

To handle 2D page layouts without needing explicit bounding-box coordinates, Qwen2.5-VL uses \textbf{2D Rotary Position Embeddings (2D-RoPE)} \cite{su2024roformer}. In 2D-RoPE, attention channels are split into separate horizontal ($x$) and vertical ($y$) rotary components. The attention between two tokens depends directly on their relative 2D distance $(\Delta x, \Delta y)$ on the page. This lets the model understand column alignment in matrices and fraction baselines directly through its attention mechanism.

\subsection{Conversational Mode Collapse}
Even with all this modeling capacity, instruction-tuned VLMs present practical headaches when used for raw transcription \cite{chen2024hallucination}:
\begin{enumerate}
    \item \textbf{Chatty Introductions:} Because the models were fine-tuned with RLHF to be helpful chatbots, they often start their output with conversational small talk (*"Here is the transcription of the handwritten notes shown in the image..."*) instead of directly outputting \LaTeX\ code.
    \item \textbf{Python Code Output:} Because their pretraining corpus includes massive amounts of GitHub code, models often treat math problems as programming prompts, generating NumPy scripts (\texttt{import numpy as np; A = np.array(...)}) instead of \LaTeX\ math environments.
\end{enumerate}

\section{Formal Languages, the Chomsky Hierarchy, and Pushdown Automata in Document Parsing}
\label{sec:formal_grammars_parsing}

Recognizing handwritten mathematics is fundamentally distinct from plain natural language processing due to the underlying formal language structure \cite{hopcroft2001introduction}. In the Chomsky Hierarchy of formal grammars, natural language text can often be approximated locally by statistical $n$-gram or finite-state regular models ($L_{\text{Reg}}$). In contrast, mathematical document markup---specifically \LaTeX---is a \textbf{Deterministic Context-Free Language (DCFL)} generated by a context-free grammar $G = (V, \Sigma, R, S)$.

\subsection{The Dyck Language of Nested Mathematical Delimiters}
Central to the syntax of \LaTeX\ is the generalized Dyck language $D_n$ defined over $n$ pairs of structured opening and closing delimiters:
\begin{equation}
\label{eq:dyck_language}
D_n = \left\{ w \in \Sigma^* \mid \text{all delimiter pairs are balanced and properly nested} \right\}
\end{equation}
where the delimiter alphabet includes environment boundaries:
\begin{equation}
\texttt{\textbackslash begin\{env\}} \iff \texttt{\textbackslash end\{env\}}, \quad \{ \iff \}, \quad [ \iff ], \quad \$ \iff \$
\end{equation}
A string $w$ belongs to $D_n$ if and only if every opened delimiter is closed in strict Last-In, First-Out (LIFO) order without cross-nesting.

\subsection{Theoretical Limitations of Pure Neural Autoregressive Decoders}
Standard autoregressive Transformers operate over a sequence of hidden vectors using soft dot-product attention:
\begin{equation}
\text{Attention}(\mathbf{Q}, \mathbf{K}, \mathbf{V}) = \text{softmax}\left( \frac{\mathbf{Q}\mathbf{K}^T}{\sqrt{d_k}} \right) \mathbf{V}
\end{equation}
While attention mechanisms can theoretically simulate bounded counter machines over short sequences, in practice they lack an explicit, unbounded stack memory. As sequence length $T$ grows ($T > 1000$ tokens), attention weights diffuse over hundreds of past positions. Consequently, pure neural decoders suffer from a non-zero probability of:
\begin{enumerate}
    \item \textbf{Scope Omission:} Failing to emit closing delimiters before end-of-sequence (e.g., generating $\texttt{\textbackslash begin\{bmatrix\}}$ but reaching \texttt{<|im\_end|>} before emitting $\texttt{\textbackslash end\{bmatrix\}}$).
    \item \textbf{Cross-Nesting Violation:} Emitting closing tokens out of order (e.g., $\texttt{\textbackslash begin\{equation\}} \dots \texttt{\textbackslash begin\{bmatrix\}} \dots \texttt{\textbackslash end\{equation\}} \dots \texttt{\textbackslash end\{bmatrix\}}$).
\end{enumerate}

Because the \TeX\ compiler operates under zero syntax tolerance---halting compilation upon encountering even a single unclosed delimiter---a purely probabilistic neural decoder cannot mathematically guarantee a $100\%$ compilation pass rate.

\subsection{The Neural-Symbolic Hybrid: Coupling VLMs with Pushdown Automata}
This theoretical insight directly motivates our architectural design: rather than attempting to force a probabilistic neural network to learn hard discrete invariants through brute-force data training, we couple the neural model with an external symbolic \textbf{Deterministic Pushdown Automaton (DPDA)} equipped with a physical LIFO stack. The neural vision-language model is assigned the task it excels at---probabilistic visual perception and token recognition ($P(\mathbf{Y} \mid \mathbf{I})$)---while the deterministic automaton guarantees syntactic closure under the Dyck grammar, ensuring $\text{Compile}(\hat{\mathbf{Y}}) = 1$ unconditionally.

\section{Summary of Our Architectural Hypotheses}
\label{sec:architectural_synthesis}

Our literature survey highlighted two fundamentally different ways to tackle handwritten math exams:
\begin{itemize}
    \item \textbf{The Modular Hybrid Approach (Divide and Conquer):} Break the task into distinct steps: clean the image, segment lines, send text to TrOCR and formulas to DenseNet HMER, and stitch everything back together with a 2D DAG. This gives total transparency into each step, but errors in early stages multiply through the rest of the pipeline.
    \item \textbf{The Multimodal VLM Approach (End-to-End):} Feed the whole page into a high-capacity vision-language model (Qwen2.5-VL) and let it parse text, math, and layout all at once. This avoids line-slicing errors, but requires strict prompt constraints to stop the model from chatting or solving the math itself.
\end{itemize}

In this project, we built and tested both paradigms on authentic exam scripts from IIT Guwahati to see how they compare in real-world transcription accuracy, GPU resource usage, and robustness against silent over-correction.
